\section{The exact weight-seven normalization}
\label{sec:weight-seven-partial}

Let $G=\widehat{\mathrm{GT}}$, let $H$ be its Galois image, and use
the established normalization
\[
 D_5=N_{\mathrm{rad}}\cap\ker m_3\cap\ker m_5,
 \qquad J_5=D_5\cap H,
 \qquad G=D_5H.
\]
All Magnus coefficients below are $2$-adic, with $x=1+S$ and $y=1+T$.
We prove that the first possible nonzero degree on $D_5$ is seven and
write $m_7(g)=[S^6T]f$ for $g=(1,f)\in D_5$.

\begin{theorem}\label{thm:weight-seven-bounds}
Every word belonging to $D_5$ has zero Magnus components below degree
seven. The map $m_7:D_5\to\mathbf Z_2$ is a continuous homomorphism and
\[
 m_7(hgh^{-1})=\lambda_2(h)^7m_7(g)
 \qquad(h\in G,\ g\in D_5).
\]
Its arithmetic and full images are exactly
\begin{equation}\label{eq:weight-seven-image-interval}
 m_7(J_5)=m_7(D_5)=128\mathbf Z_2.
\end{equation}
The corresponding degree-seven Hall images both equal
$8\mathbf Z_2\mathbf L$, where $\mathbf L$ is the primitive
integral vector specified below.
\end{theorem}

\subsection{The first possible leading term}

Use standard Lyndon brackets on $x<y$, ordered first by degree and
then lexicographically, and let $h_w$ denote the corresponding integral
Hall coordinate. The standard factorization $w=uv$ uses the longest
proper Lyndon suffix, and the group word is $L_w=[L_u,L_v]$.

\begin{lemma}\label{lem:weight-seven-leading-spaces}
Suppose that $f\in\gamma_5$ satisfies the two-cycle and three-cycle
relations with multiplier one. Put $a=[S^4T]f$. Its Hall coordinates
in degrees five and six are
\begin{equation}\label{eq:weight-seven-lower-vector}
 a\left(1,-2,\tfrac32,2,\tfrac12,-1,
 -2,\tfrac72,-\tfrac32,-4,-\tfrac12,\tfrac32,
 \tfrac72,\tfrac32,-2\right).
\end{equation}
In particular, $a=0$ implies that $f$ starts in degree seven.

If $f$ starts in degree seven and satisfies the same two relations,
its degree-seven Hall vector has the form $a\mathbf v+b\mathbf w$,
where $a=[S^6T]f$, $b$ is the sixth degree-seven Hall coordinate, and
\begin{align}
\mathbf v={}&\left(1,-3,5,5,-2,0,-12,-5,-2,\tfrac23,
                  -5,5,8,3,-\tfrac{11}3,4,3,-1\right),
                  \label{eq:weight-seven-v}\\
\mathbf w={}&\left(0,0,0,0,0,1,1,0,0,\tfrac13,
                  -1,-1,-1,0,-\tfrac13,0,0,0\right).
                  \label{eq:weight-seven-w}
\end{align}
These are necessary leading-term conditions; they assert no lifting
of an arbitrary truncated solution to a profinite automorphism.
\end{lemma}

\begin{proof}
The finite identities in Section~\ref{subsec:weight-seven-leading-certificate}
give \eqref{eq:weight-seven-lower-vector} and
\eqref{eq:weight-seven-v}--\eqref{eq:weight-seven-w} directly from
coefficients of the two defining relations through degree nine.
They are exact rational identities between explicitly specified
integral coefficient forms. Consequently they apply to the
$2$-adic Hall coordinates of a word.

On $D_5$, the preceding cubic calculation first gives $f\in\gamma_5$,
and $m_5=0$. Equation~\eqref{eq:weight-seven-lower-vector} therefore
makes every degree-five and degree-six coordinate zero. This proves
the asserted absence of lower terms on $D_5$.
\end{proof}

\subsection{The additional leading pentagon relation}

The original pure-braid pentagon has the form
\[
 f(x_{12},x_{23}x_{24})f(x_{13}x_{23},x_{34})
 =f(x_{23},x_{34})f(x_{12}x_{13},x_{24}x_{34})f(x_{12},x_{23});
\]
see \cite[Remark~1.2, p.~98]{Furusho2003}.
For a word starting in degree seven, its degree-seven term $p$
therefore satisfies
\begin{equation}\label{eq:weight-seven-leading-pentagon}
\begin{split}
 p(t_{12},t_{23}+t_{24})+p(t_{13}+t_{23},t_{34})
 -p(t_{23},t_{34})\hspace{12mm}\\
 -p(t_{12}+t_{13},t_{24}+t_{34})-p(t_{12},t_{23})=0.
\end{split}
\end{equation}

\begin{lemma}\label{lem:weight-seven-integral-line}
For an element of $D_5$, the parameters in
\eqref{eq:weight-seven-v}--\eqref{eq:weight-seven-w} satisfy
\[
 16b=19a.
\]
Consequently $a\in16\mathbf Z_2$, and the integral leading line is
generated by the primitive vector
\begin{equation}\label{eq:weight-seven-primitive-vector}
\mathbf L=(16,-48,80,80,-32,19,-173,-80,-32,
           17,-99,61,109,48,-65,64,48,-16).
\end{equation}
The leading Hall vector of an element is $(a/16)\mathbf L$.
Moreover, its leading alternating Magnus coefficient is
$[STSTSTS]f=-2a$.
\end{lemma}

\begin{proof}
Evaluate \eqref{eq:weight-seven-leading-pentagon} in the infinitesimal
Artin action on the free Lie algebra on $X_1,\ldots,X_4$. Write
\[
 D_{ij}(X_i)=[X_i,X_j],\qquad
 D_{ij}(X_j)=[X_j,X_i],\qquad D_{ij}(X_k)=0\quad(k\ne i,j).
\]
The relations $[D_{ij},D_{kl}]=0$ for four distinct indices and
$[D_{ij},D_{ik}+D_{jk}]=0$ follow from Jacobi. A common change of sign
in this infinitesimal action does not change the required vanishing
in degree seven. No faithfulness assertion is needed.

Take the coefficient of $X_1^4X_2X_3X_2X_4$ in the image of $X_1$.
For $p=a\mathbf v+b\mathbf w$, the five terms, before attaching
their respective signs $+,+,-,-,-$, give
\[
\begin{array}{c|rr}
\text{term}&\text{coefficient of }a&\text{coefficient of }b\\ \hline
1&46/3&-16/3\\
2&-6&0\\
3&0&0\\
4&3&0\\
5&0&0.
\end{array}
\]
This table follows by the displayed formulas for $D_{ij}$ and the
recursion $[D,E](X_i)=D(E(X_i))-E(D(X_i))$, applied to the standard
Lyndon brackets. The pentagon thus gives $(19a-16b)/3=0$.

Both $a$ and $b$ are integral Hall coordinates, and $19$ is odd.
Hence $a\in16\mathbf Z_2$. Substituting $b=19a/16$ into
\eqref{eq:weight-seven-v}--\eqref{eq:weight-seven-w} gives
\eqref{eq:weight-seven-primitive-vector}; its coordinate $19$ proves
primitivity. Direct bracket expansion on the preceding plane gives
$[STSTSTS]p=(32/3)(a-b)$, which becomes $-2a$.
\end{proof}

\subsection{The exact beta coordinate and an arithmetic value}

For a multiplier-one element, put $a_j=[S^jT]f$ and define
\begin{equation}\label{eq:weight-seven-tau}
 \tau_7(g)=360a_6+900a_5+780a_4-239a_2\in\mathbf Z_2.
\end{equation}

\begin{lemma}\label{lem:weight-seven-beta}
The map $\tau_7$ is a continuous homomorphism on $\ker\lambda$, with
conjugation weight $\lambda_2^7$. On $D_5$ it equals $360m_7$.
\end{lemma}

\begin{proof}
Use the normalized Alexander beta cofactor
\[
 B_g(X,Y)=1+(X-1)(Y-1)R_g(X,Y),\qquad [f]=R_g(X,Y)[x,y].
\]
The action on the rank-one Alexander module has scalar
$S_\lambda(X)S_\lambda(Y)B_g(X,Y)$, where
$S_\lambda(X)=(X^\lambda-1)/(X-1)$. The $S_\lambda$ factors satisfy
the same semilinear multiplicative cocycle identity as that scalar;
cancelling them gives the corresponding identity for $B_g$.
Thus
\[
 \vartheta_7(g)=[u^6v]\log B_g(e^u,e^v)
\]
is an additive $\mathbf Q_2$-valued homomorphism on $\ker\lambda$,
with conjugation weight $\lambda_2^7$.

Since $B_g(X,1)=1$, differentiation at $v=0$ gives
\[
 \left.\partial_v\log B_g(e^u,e^v)\right|_{v=0}
   =(e^u-1)R_g(e^u,1)
   =\sum_{j\ge2}a_j(e^u-1)^j.
\]
The last equality follows from the affine specialization
$x=\operatorname{diag}(1+S,1)$, $y=I+TE_{12}$: the commutator has
translation $ST$, conjugation by $x$ multiplies it by $1+S$, and
conjugation by $y$ acts trivially on that translation module.
The coefficients of $u^6$ in $(e^u-1)^j$, for $j=2,\ldots,6$, are
\[
 \frac{31}{360},\quad\frac34,\quad\frac{13}{6},\quad\frac52,
 \quad1.
\]
Using the earlier identity $a_3=-a_2$ therefore gives
\[
 \vartheta_7=a_6+\frac52a_5+\frac{13}{6}a_4
                         -\frac{239}{360}a_2.
\]
Multiplication by $360$ proves the integral formula
\eqref{eq:weight-seven-tau}, continuity and the exact weight law.
On $D_5$, all the terms except $a_6$ vanish by
Lemma~\ref{lem:weight-seven-leading-spaces}.
\end{proof}

\begin{proposition}\label{prop:weight-seven-arithmetic-value}
There is $z_5\in J_5$ with
\[
 m_7(z_5)\in128\mathbf Z_2^\times.
\]
More explicitly, it can be chosen so that
$m_7(z_5)=128\cdot199017\,\tau_7(\sigma)$ for an arithmetic
$\sigma$ with $\tau_7(\sigma)$ odd.
\end{proposition}

\begin{proof}
Simultaneous rational Kummer surjectivity supplies
$\sigma\in H\cap\ker\lambda$ with $\rho_2^I(\sigma)=2$ in
$\widehat{\mathbf Z}$ and every odd-prime rational Kummer coordinate
zero. The established boundary congruence
$\rho_2^I\equiv-2m_3\pmod4$ makes $a_2(\sigma)$ odd.
Equation~\eqref{eq:weight-seven-tau} consequently makes
$\tau_7(\sigma)$ odd.

Choose $h\in H$ with $\lambda_2(h)=3$ and put
$t=\lambda(h)\in\widehat{\mathbf Z}^{\times}$. For fixed $h$, set
\[
 \mathcal D_c(z)=hzh^{-1}z^{-c},\qquad
 z_1=\mathcal D_t(\sigma),\quad
 z_3=\mathcal D_{t^3}(z_1),\quad
 z_5=\mathcal D_{t^5}(z_3).
\]
The powers are the continuous profinite powers in the procyclic
closure of the element. The first operation kills every rational
Kummer coordinate. The second kills $m_3$ and places $z_3$ in $J_3$.
The third kills $m_5$ by its established weight-five transformation
law on $D_3$. All preceding vanishings are preserved, so $z_5\in J_5$.

The exact homomorphism and weight law of $\tau_7$ give
\[
 \tau_7(z_5)=(3^7-3)(3^7-3^3)(3^7-3^5)\tau_7(\sigma).
\]
Since $\tau_7=360m_7$ on $D_5$, division by $360$ gives
\[
 m_7(z_5)=\frac{2184\cdot2160\cdot1944}{360}\tau_7(\sigma)
         =128\cdot199017\,\tau_7(\sigma).
\]
The last two factors are odd. The compact arithmetic image
therefore contains $128\mathbf Z_2$; the uniform odd theorem
gives the reverse inclusion.
\end{proof}

\begin{proof}[Proof of Theorem~\ref{thm:weight-seven-bounds}]
The lower-degree vanishing and integral leading line were proved in
Lemmas~\ref{lem:weight-seven-leading-spaces} and
\ref{lem:weight-seven-integral-line}. The beta-coordinate identity
gives additivity and the scalar conjugation rule for $m_7$; division
by the nonzero integer $360$ is legitimate in the torsion-free group
$\mathbf Z_2$.

The arithmetic value and closedness show that both images contain
$128\mathbf Z_2$. Since $D_5=C_2$,
Theorem~\ref{thm:uniform-odd-arithmetic-image} with $r=3$ gives
the reverse inclusion. Since the leading vector is
$(m_7/16)\mathbf L$, both Hall images are exactly
$8\mathbf Z_2\mathbf L$.
\end{proof}

\subsection{Finite identities for the leading spaces}
\label{subsec:weight-seven-leading-certificate}

Here are finite recursions and explicit matrix identities that verify
Lemma~\ref{lem:weight-seven-leading-spaces}. They specify the
coefficient calculation without a dependence on an implementation.
Work in noncommutative integer polynomials truncated after degree
nine. Put $X=1+S$, $Y=1+T$, $Z=Y^{-1}X^{-1}$ and
$(1+A)^{-1}=\sum_{j=0}^{9}(-A)^j$. Form the standard Lyndon group
words by the preceding commutator recursion.

For $d=5$ or $7$, let the columns be the Lyndon words of lengths
$d,\ldots,9$, ordered first by length and then lexicographically.
There are $119$ columns for $d=5$ and $104$ for $d=7$. Set
\[
 F_d=1+\sum_{d\le |w|\le9}h_w(L_w-1).
\]
All coefficients through degree nine are linear in the $h_w$,
because products of two nonconstant factors start in degree at
least ten. Define the rows of the integer coefficient matrix $A_d$
by
\[
 I_v=[v]\bigl(F_d(Y,X)+F_d(X,Y)-2\bigr),
\]
\[
 II_v=[v]\bigl(F_d(Z,X)+F_d(Y,Z)+F_d(X,Y)-3\bigr).
\]
The rows vanish on the coordinates of any word satisfying
the two-cycle and three-cycle relations.

Encode a word $v$ of length $n$ by
$c(v)=2^n+\operatorname{bin}(v)$, where $S$ is the binary digit $0$
and $T$ is $1$. The row codes for $I_v$ and $II_v$ are respectively
$c(v)$ and $1024+c(v)$. Let $R_5,R_7$ be the row lists displayed
below, let $C_5=(1,\ldots,96)$ and $C_7=(1,\ldots,82)$, and put
$B_d=A_d[R_d,C_d]$.

Let $\mathbf u$ be the fifteen-entry vector in
\eqref{eq:weight-seven-lower-vector}. In the $119$-dimensional
row space let $T_5$ have the fourteen rows
$2(e_j-u_je_1)$, $2\le j\le15$. In the $104$-dimensional row
space let $T_7$ have the sixteen rows
$3(e_j-v_je_1-w_je_6)$ for $1\le j\le18$, $j\ne1,6$.
Both matrices have integral entries. The finite identities are
\begin{equation}\label{eq:weight-seven-finite-matrix-identity}
 \det B_5=-6603882696,\qquad \det B_7=42661728,
 \qquad
 T_d[:,C_d]B_d^{-1}A_d[R_d,:]=T_d\quad(d=5,7).
\end{equation}
Every entry in these identities is specified by the integer
recursions and the row lists. The inverses are taken over
$\mathbf Q$, and the identities are exact, not congruences.
Applying them to $A_dh=0$ gives $T_dh=0$ and hence the two leading
vector formulas. This is a proof of necessary vanishings and
leading relations; it supplies no lifting assertion for the
remaining coordinates.

% The explicit row lists are appended immediately below.

\[
\begin{split}
R_{5}=\bigl( &33,35,37,65,67,69,77,129,131,133,135,137,\\
 &139,141,147,149,257,259,261,263,265,267,269,275,\\
 &277,281,283,285,293,301,513,515,517,519,521,523,\\
 &525,527,529,531,533,535,537,539,541,547,549,551,\\
 &553,555,557,563,565,569,587,589,595,597,1059,1089,\\
 &1093,1099,1153,1155,1157,1161,1163,1165,1281,1283,1285,1287,\\
 &1289,1291,1299,1301,1317,1325,1537,1539,1541,1545,1547,1549,\\
 &1553,1555,1557,1559,1561,1563,1571,1573,1577,1579,1581,1589\bigr).
\end{split}
\]

\[
\begin{split}
R_{7}=\bigl( &129,131,133,135,137,139,141,147,149,257,259,261,\\
 &263,265,267,269,275,277,281,283,285,293,301,513,\\
 &515,517,519,521,523,525,527,529,531,533,535,537,\\
 &539,541,547,549,551,553,555,557,563,565,569,587,\\
 &589,595,597,1155,1157,1161,1281,1283,1285,1287,1289,1291,\\
 &1299,1301,1317,1325,1537,1539,1541,1545,1547,1549,1553,1555,\\
 &1557,1559,1561,1563,1571,1573,1577,1579,1581,1589\bigr).
\end{split}
\]
